\documentclass[11pt,a4paper]{article}
\usepackage[top=2.5cm,bottom=2.8cm,left=2.5cm,right=2.5cm]{geometry}
\usepackage{amsmath,amssymb,amsthm}
\usepackage{hyperref,xcolor,enumitem,booktabs,array}
\usepackage{graphicx,subcaption}

\hypersetup{colorlinks=true,linkcolor={blue!60!black},
  citecolor={blue!60!black},urlcolor={blue!60!black}}

\newtheorem{theorem}{Theorem}[section]
\newtheorem{lemma}[theorem]{Lemma}
\newtheorem{proposition}[theorem]{Proposition}
\newtheorem{corollary}[theorem]{Corollary}
\theoremstyle{definition}\newtheorem{definition}[theorem]{Definition}
\theoremstyle{remark}
\newtheorem{observation}[theorem]{Observation}
\newtheorem{remark}[theorem]{Remark}
\newtheorem{conjecture}[theorem]{Conjecture}

\newcommand{\vp}{\varphi}
\newcommand{\TCMB}{T_{\mathrm{CMB}}}
\newcommand{\Lcond}{\Lambda_{\mathrm{cond}}}
\newcommand{\rCMB}{\rho_{\mathrm{CMB}}}
\newcommand{\SU}{\mathrm{SU}}
\newcommand{\twoI}{2I}
\newcommand{\twoT}{2T}
\newcommand{\QH}{Q_H}
\newcommand{\Qgr}{Q_{\mathrm{group}}}
\newcommand{\Ztwo}{\mathbb{Z}_2}
\newcommand{\Zthree}{\mathbb{Z}_3}
\newcommand{\Tw}{\mathrm{Tw}}
\newcommand{\Mup}{M_{\mathrm{up}}}
\newcommand{\Mdn}{M_{\mathrm{down}}}
\newcommand{\alphas}{\alpha_s}
\newcommand{\alphaem}{\alpha_{\mathrm{em}}}

\title{\textbf{The Dynamical Layer of the Quark Sector:\\
$E_6$-Native Generations, the Isospin Mass-Scale, and\\
Static Residuals as Strong Dynamics}
\\[0.6em]
{\large Companion Paper XX to the Density-Feedback
Faddeev--Niemi Hopfion Series}
\\[0.6em]
\large DRAFT --- September 2026}
\author{Frederick Manfredi}
\date{}

\begin{document}
\maketitle

\begin{abstract}
Papers~V, XV--XIX establish the $\QH=3$ quark/baryon sector of the
density-feedback Faddeev--Niemi Hopfion: the trefoil $T_{2,3}$, the
binary tetrahedral group $\twoT$, the WZW model $(E_6)_1\supset
\SU(3)_1^{\times3}$, and an algebraic (saddle-less) quark mass. This
paper assembles the sector's \emph{dynamical} content and, in doing so,
reframes what the static/topological framework is and is not.
We record six results and one organising principle.
(1)~\textbf{$E_6$-native generations, $E_8$-bridge anchor.} The quark's
own group is $E_6$; $E_8$ enters only as the charged-lepton bridge that
anchors the down sector (strange$=$muon). The generation ladder is the
per-strand $E_6$ ribbon twist; the up/down tower-steepness asymmetry is
the ratio of WZW phase units $\pi\Qgr/(k\,h)$, namely
$(\pi/4)/(\pi/9)=9/4$, matching the measured steepness ($\approx2.1$)
to $6\%$ with no free parameter.
(2)~\textbf{Isospin mass-scale from tangent orientation.} The up/down
mass-scale ratio is fixed by the trefoil crossing/midpoint tangent
geometry (Paper~XVIII), $\Mdn/\Mup\approx0.189$, closing to the measured
$0.157$ under a single generation-universal correction; a direct
Faddeev-energy computation shows the operative quantity is tangent
\emph{orientation}, not a static energy difference.
(3)~\textbf{Fractional charge from colour triality.} The $\Zthree$
center of the colour CFT $\SU(3)_1=(E_6)_1$ forces the quark hypercharge
$Y\equiv\text{triality}/3$, giving $Q=\{+\tfrac23,-\tfrac13\}$ for the
colour triplet and integer charge for the colour-singlet leptons; the
fractional denominator is the colour center (the three crossings of
$T_{2,3}$), and the Chern--Simons edge modes are its dynamical carrier.
This derives what Paper~XVIII assigned by writhe.
(4)~\textbf{The generation-1 electromagnetic flip.} A charge-linear
coupling of scale $\alphaem\Lambda_{\mathrm{const}}\sim$~MeV flips the
$u<d$ ordering at generation~1 while leaving $c>s$, $t>b$ intact.
(5)~\textbf{The doublet sum rule, and the constituent scale.}
$m_{\rm up}+m_{\rm down}=C\,m_\ell$ with $C\approx13$ holds across the two
confined doublets and cross-predicts the confined up quarks. The
constituent scale is derived, $m_{\rm const}=\Lcond^{\rm baryon}\vp^6
e^{8\pi}\approx237$~MeV (the $8\pi$ Chern--Simons spoke), and reduces $C$
to the $E_6$-native tower level: deriving $C$ and pinning that tower are
one problem, with the tower level the sole remaining unknown.
(6)~\textbf{Residuals as the dynamical layer.} A diagnostic across all
sectors shows the static/topological predictions are exact at leading
order for static-saddle sectors (charged leptons, gauge couplings:
residuals $\lesssim0.5\%$) and leave $\sim\!\alphas$ residuals
($\sim20$--$30\%$) precisely in the dynamical sectors (quarks, neutrino).
The residuals are not unrelated misses but one correction whose size
tracks the sector's coupling, largest where the object has no static
saddle. We record, as retracted, several numerological readings of $C$
(the $4\pi$, $\vp^6$, and solid-angle interpretations), and identify the
Chern--Simons edge/Hall dynamics and the non-perturbative formation
energy as the two un-developed dynamical handles.
\end{abstract}

\tableofcontents
\bigskip\newpage

%══════════════════════════════════════════════════════════════════════════════
\section{Introduction}
\label{P20:sec:intro}
%══════════════════════════════════════════════════════════════════════════════

The density-feedback Faddeev--Niemi Hopfion framework~\cite{Paper1}
assigns each topological sector $\QH=N$ an identity through a
McKay-correspondence group, a WZW conformal field theory, and a knot
type (Paper~XVII~\cite{Paper17}). For the $\QH=3$ quark/baryon sector this
triple is $\bigl(\twoT,\ (E_6)_1\supset\SU(3)_1^{\times3},\ T_{2,3}\bigr)$
(Papers~XV, XVI~\cite{Paper15,Paper16}). The lepton sectors are
\emph{static}: the charged lepton is a stable Faddeev saddle whose mass
is the configuration energy, and the framework predicts it to
$1.3\times10^{-4}$ (Paper~IV~\cite{Paper4}, Paper~XII~\cite{Paper12}).
The quark sector is not: the isolated trefoil is metastable with no
stable field saddle (Paper~XV~\cite{Paper15}, Paper~XVIII~\cite{Paper18}),
so the quark mass is algebraic rather than a saddle energy
(Paper~V~\cite{Paper5}), and the framework's quark-mass predictions
carry residuals an order of magnitude larger than the lepton sector's.

This paper collects the sector's dynamical structure. Sections~\ref{P20:sec:groups}--\ref{P20:sec:gen1}
establish the group and geometric content of the generation, isospin,
and charge structure (including a derivation of the fractional charge
from the colour $\Zthree$ center, Section~\ref{P20:sec:charge});
Section~\ref{P20:sec:sumrule} records an empirical doublet
mass-budget relation and, importantly, what it does \emph{not} mean;
Section~\ref{P20:sec:residuals} presents the organising result---a
diagnostic showing that the framework's residuals partition cleanly by
whether the sector is a static saddle or a dynamical object, with the
dynamical residuals sitting at the strong-coupling scale. The
consistent theme is that the static/topological layer is the exact
leading order and the residuals are a single dynamical correction,
forced into view in the quark sector because the quark is intrinsically
dynamical.

We separate established results (theorems, propositions) from leads
(conjectures) and from honest negatives (remarks). Several attractive
numerical readings are recorded as \emph{retracted} to prevent their
recurrence.

Several of the structures used here have established counterparts, which
we cite rather than claim: the golden ratio $\vp$ is the quantum
dimension of the Fibonacci anyon of the $\SU(2)_3$ WZW
model~\cite{Nayak2008}; icosahedral ($A_5$) family and modular flavour
symmetry produce golden-ratio lepton mixing and hierarchical fermion
masses from the modulus fixed points~\cite{Everett2009,Feruglio2017,ModFlavHierarchy2025};
and the fractional-charge quantisation of Section~\ref{P20:sec:charge}
is the standard $\mathbb{Z}_6$ quotient of the Standard-Model gauge
group~\cite{FracCharge2024}. The framework's contribution is not these
individual facts but a single topological origin---the density-feedback
hopfion with its McKay/WZW data---from which they follow together,
alongside the mass and residual results below; in particular the
mass-hierarchy-from-fixed-point mechanism of modular flavour symmetry is
paralleled by the RG-fixed-point tower ($\zeta=\vp$, Paper~VI~\cite{Paper6}).

%══════════════════════════════════════════════════════════════════════════════
\section{Generation Structure: $E_6$-Native, $E_8$-Bridge}
\label{P20:sec:groups}
%══════════════════════════════════════════════════════════════════════════════

\subsection{The two groups and their roles}

Paper~V~\cite{Paper5} first obtained the quark masses from an $E_8$
Coxeter route applied to the entire scale. Papers~XVI, XVII refine this:
the quark's intrinsic group is $E_6$ (the McKay dual of the $\vp$-free
group $\twoT$, with $(E_6)_1\supset\SU(3)_1^{\times3}$ the colour CFT),
while $E_8$ (the McKay dual of $\twoI$, the charged-lepton group) enters
only as a \emph{bridge}: the down sector is anchored to the charged
leptons through the exact coincidence $m_s=m_\mu$ at the condensate
scale (the strange--muon anchor, forced by $T^{(13)}=\tfrac18=T_\mu$;
Paper~XVI~\cite{Paper16}). %\ref{P17:rem:quark_mass_algebraic}

\begin{definition}[Native group and bridge]
\label{P20:def:native_bridge}
For the $\QH=3$ sector we call $E_6$ the \emph{native} group (it
carries the colour CFT $\SU(3)_1$ and the per-strand generation twist)
and $E_8$ the \emph{bridge} (it references the charged-lepton mass that
anchors the down-type quarks). The down-type masses are computed on the
$E_8$ bridge; the up-type generation structure is $E_6$-native.
\end{definition}

\subsection{The up/down steepness as a ratio of WZW phase units}

The mass formula of Paper~VI~\cite{Paper6}, Paper~XVII~\cite{Paper17}
carries a WZW phase unit $\pi\Qgr/(k\,h)$, where $\Qgr$ is the group
order, $k$ the WZW level and $h$ the Coxeter number. For the charged
lepton (the $E_8$ bridge), $\Qgr=10$, $k=3$, $h(E_8)=30$, giving the
established unit $\pi/9$. The $E_6$-native quark carries its own numbers
$\Qgr=3$, $k=1$ ($\SU(3)_1=(E_6)_1$), $h(E_6)=12$.

\begin{proposition}[Up/down tower-steepness ratio]
\label{P20:prop:steepness}
The $E_6$-native phase unit is
\begin{equation}
  \frac{\pi\,\Qgr}{k\,h(E_6)}=\frac{\pi\cdot3}{1\cdot12}=\frac{\pi}{4},
  \label{P20:eq:e6phase}
\end{equation}
and its ratio to the $E_8$-bridge unit $\pi/9$ is
$(\pi/4)/(\pi/9)=9/4=2.25$. The up-type tower is therefore steeper than
the down-type by this factor. The empirically observed ratio of
per-generation tower increments, $\Delta n_{\rm up}/\Delta n_{\rm down}
\approx2.1$ (confined generations), agrees to $6\%$ with no free
parameter.
\end{proposition}

\begin{remark}[Generation is a ribbon twist, not a curve property]
\label{P20:rem:twist}
Generation is realised as the per-strand $E_6$ ribbon twist $|\Tw|\in
\{1,2,5\}$ (the Coxeter-exponent distances $|m-h/2|$ for the three
mirror-pair doublets), a framing of the fixed knot $T_{2,3}$; the knot
type, writhe (colour), and confinement length are generation-universal.
The mirror symmetry $\Tw\to-\Tw$ is the isospin doublet, and the
framing period is $\mathrm{mod}\ 3$ (the $\SU(3)_1$ topological spin),
consistent with quarks carrying colour.
\end{remark}

\begin{remark}[The empirical down power law is not a fundamental identity]
\label{P20:rem:powerlaw}
The confined down masses obey a clean power law $m\propto|\Tw|^{p}$ with
$p\approx4.2$ ($R^2=0.9999$). This is \emph{not} a representation-theory
identity: the per-step exponent scatters ($4.17$ vs $4.22$), and the
clean power requires the measured charged-lepton masses as inputs (the
WZW data alone give $4.03$ vs $4.57$). It is a re-parametrisation of the
$E_8$-bridge route (down mass $=m_\ell$ times the $E_8$-exponent
factor), not a new law; the proximity of $p$ to $\vp^3=4.236$ lies
within the fit scatter and is not claimed.
\end{remark}

\begin{proposition}[Parameter-free up-doublet ratio]
\label{P20:prop:up_ratio}
Write the $E_6$-native tower level as $n_g=A\,T_g-2\,m_g^{E_6}$, with
$A=6\,h(E_6)=72$ (the $E_6$ analogue of the $E_8$-bridge coefficient
$6\,h(E_8)=180$; Paper~XVII~\cite{Paper17}), $T_g$ the $\SU(2)_{k_g}$
T-matrix phases $\{5/24,\,1/8,\,3/40\}$ at the spinor primary $j=\tfrac12$,
$k_g=1,2,3$ (Paper~XVII~\cite{Paper17}), and $m_g^{E_6}$ the $E_6$ Coxeter
exponents of the up doublet, $\{m_u,m_c,m_t\}=\{7,8,11\}$
(Remark~\ref{P20:rem:twist}). The first-to-second generation gap is an
integer,
\begin{equation}
  n_u-n_c=A\,(T_1-T_2)-2\,(m_u^{E_6}-m_c^{E_6})
        =72\cdot\tfrac{1}{12}-2\,(7-8)=6+2=8,
  \label{P20:eq:up_gap}
\end{equation}
so with the $E_6$-native phase unit $\pi/4$~\eqref{P20:eq:e6phase} and the
tower $m_g\propto\exp(-n_g\,\pi/4)$,
\begin{equation}
  \frac{m_u}{m_c}=\exp\!\Bigl[-(n_u-n_c)\tfrac{\pi}{4}\Bigr]
                =e^{-2\pi}=1.867\times10^{-3},
  \label{P20:eq:up_ratio}
\end{equation}
a mass ratio fixed entirely by WZW T-matrix and Coxeter data, with no free
parameter---the absolute offset common to all $n_g$ cancels in the ratio.
\end{proposition}

\begin{remark}[Empirical status and the open anchor]
\label{P20:rem:up_ratio_data}
The QCD mass anomalous dimension is flavour-independent, so the ratio
$m_u/m_c$ evaluated at a common renormalisation scale is scale-invariant.
Running $m_c(m_c)=1.27$~GeV up to $2$~GeV gives $m_c=1.10$~GeV, whence
$m_u/m_c=2.16/1097=1.97\times10^{-3}$ at a common scale (identical at $M_Z$),
against the raw mixed-scale value $2.16/1270=1.70\times10^{-3}$. The
prediction $e^{-2\pi}=1.87\times10^{-3}$ lies between the two and well within
the dominant uncertainty, that of $m_u$ ($\sim\pm20\%$). This extends the
``conformal weight $=$ observable'' identity of Paper~XV~\cite{Paper15}
---there established for the electric charges---from a charge to a mass ratio.
The ratio fixes only the tower \emph{spacing}; the absolute normalisation is
not determined by it and is the content of Open Problem~O1. Crucially, the top
is \emph{not} a rung of this confined tower: it is the unique quark that decays
before it can hadronise ($\tau_t/\tau_{\rm QCD}\approx0.15$), so it never
completes the $T_{2,3}$ formation and sits at $v/\sqrt2$ rather than at its
full-formation value ($m_\tau\,e^{15\pi/9}\approx334$~GeV on the $E_8$ bridge)
---an incomplete-formation deficit of $\approx2$, not a tower offset.
Extending the confined
$u,c$ tower onto the bare top is therefore inappropriate (the up triplet is
inhomogeneous: $u,c$ confined, $t$ bare). The absolute scale of the confined
$u,c$ doublet is set instead by the $T_{2,2}\!\to\!T_{2,3}$ formation energy.
\end{remark}

%══════════════════════════════════════════════════════════════════════════════
\section{The Isospin Mass-Scale from Tangent Orientation}
\label{P20:sec:isospin}
%══════════════════════════════════════════════════════════════════════════════

Paper~XVIII assigns isospin to the $\Ztwo$ writhe asymmetry of $T_{2,3}$:
the crossing-vertex network ($z=+r_0$) is up-type, the midpoint/distal
network ($z=0$) is down-type. %\ref{P18:conj:isospin}
The associated mass-scale ratio (Paper~XVIII) is set by the out-of-plane
tangent fraction at each site network.

\begin{proposition}[Isospin mass-scale and its residual]
\label{P20:prop:isospin_scale}
With the exact golden-trefoil tangent data ($R_0=3$, $r_0=\sqrt2/\vp$;
$\dot\Gamma_z=+0.321$ distal, $-0.525$ midpoint, $0$ at crossings), the
out-of-plane fraction averages $\langle\dot\Gamma_z^2\rangle=0.189$, and
the mass-scale ratio (geometric-mean quark masses per triplet) is
predicted
\begin{equation}
  \frac{\Mdn}{\Mup}
  =\frac{\langle\dot\Gamma_z^2\rangle}{1-\dot\Gamma_z^2|_{\rm crossing}}
  =0.189,
  \label{P20:eq:isospin_ratio}
\end{equation}
against the measured $0.157$. The residual is a factor $\approx1.2$,
matching in magnitude the uniform correction by which the
charged-lepton-anchored quark-mass predictions of
Paper~XVI~\cite{Paper16} overshoot; the mass-scale ratio is therefore
the tangent geometry up to a single generation-universal factor.
\end{proposition}

\begin{remark}[The operative quantity is orientation, not energy]
\label{P20:rem:orientation}
A direct evaluation of the Faddeev energy density on the crossing versus
midpoint networks (using the validated Construction-C trefoil ansatz)
gives equal densities per unit arc (quartic ratio $0.94$, quadratic
$1.00$; the energy is $z\to-z$ symmetric although the tangent
orientation is not). Substituting the local curvature, the integrated
bending $\int\kappa^2\,ds$, or the inter-strand distance for the tangent
fraction reproduces the ratio less well or with the wrong sign. The
crossing network is distinguished by its purely in-plane exit
($\dot\Gamma_z=0$), \emph{not} by any spatial proximity of its strands
or by a static energy excess. The up/down asymmetry is thus orientational
and field-selection-driven, consistent with Paper~XVIII's premise that
absent an external field the two isospin values are degenerate.
\end{remark}

%══════════════════════════════════════════════════════════════════════════════
\section{Fractional Charge from the Colour $\Zthree$ Center}
\label{P20:sec:charge}
%══════════════════════════════════════════════════════════════════════════════

The isospin assignment of Section~\ref{P20:sec:isospin} fixes
$I_3=\pm\tfrac12$ from the trefoil $z$-network; the electric charge
$Q=I_3+Y/2$ requires the hypercharge $Y$. We show the fractional part of
$Y$ is forced by the $\Zthree$ center of the colour CFT
$\SU(3)_1=(E_6)_1$, replacing the static charge assignment of
Paper~XVIII by a derivation from the colour representation.

\begin{proposition}[Fractional charge from colour triality]
\label{P20:prop:frac_charge}
The colour CFT $\SU(3)_1$ has three primaries---$\mathbf{1}$
(triality $t=0$), $\mathbf{3}$ ($t=1$), $\overline{\mathbf{3}}$
($t=2$)---with conformal weights $h=0,\tfrac13,\tfrac13$ and central
charge $c=2$. The fundamental's topological spin
$e^{2\pi i h}=e^{2\pi i/3}$ equals the $\Zthree$ center element, so
$h=\tfrac13$ encodes triality one. Consistency of the combined
$\SU(3)_{\rm colour}\otimes\mathrm{U}(1)_{\rm em}$ representation under
the common $\Zthree$ center forces the hypercharge to satisfy
\begin{equation}
  Y\equiv \frac{t}{3}\pmod 1,
  \label{P20:eq:Y_triality}
\end{equation}
with $t$ the colour triality. The quark, in the colour $\mathbf{3}$
($t=1$), therefore has $Y=\tfrac13$ and, with the $z$-network isospin
$I_3=\pm\tfrac12$ (Section~\ref{P20:sec:isospin}),
\begin{equation}
  Q=I_3+\tfrac12 Y
   =\Bigl\{+\tfrac23\ (\text{up}),\ -\tfrac13\ (\text{down})\Bigr\}.
\end{equation}
The lepton, a colour singlet ($t=0$), has integer $Y$ and hence integer
charge ($0$ neutrino, $-1$ charged). Fractional charge is thus the
signature of colour triality; the denominator $3$ is the $\Zthree$
center of $\SU(3)_1$, i.e.\ the three crossings of $T_{2,3}$.
\end{proposition}

\begin{remark}[The denominator is the colour center, not the WZW level; the Chern--Simons realisation]
\label{P20:rem:charge_cs}
Two points. First, the fractional ``$3$'' is the $\Zthree$ center of the
\emph{colour} group $\SU(3)_1=(E_6)_1$ (triality, the three crossings),
not the level $k=3$ of the charged-lepton $\SU(2)_3$; these are distinct
threes, and the charge denominator is the former. Second,
Proposition~\ref{P20:prop:frac_charge} is the group-theoretic content of
the Chern--Simons edge picture: the $\SU(3)_1$ Chern--Simons theory on a
bounded domain carries chiral edge modes which, by the bulk--edge
correspondence, carry the $\Zthree$ (triality) charge; the
fractionally-charged quark excitations are these edge modes, and the
electromagnetic Hall response fractionalises charge into thirds. This
upgrades the static writhe assignment of Paper~XVIII to a derivation
from the colour CFT. %\ref{P18:conj:isospin}
The hypercharge quantisation~\eqref{P20:eq:Y_triality} is the standard
$\mathbb{Z}_6$-quotient condition of the Standard-Model gauge group
(the combined $\SU(3)\times\SU(2)\times\mathrm{U}(1)$
centers)~\cite{FracCharge2024}; the framework-native input is that
$\SU(3)_1$ itself is derived from the trefoil $T_{2,3}$ and the binary
tetrahedral group $\twoT$
(Papers~XV--XVII~\cite{Paper15,Paper16,Paper17}).
\end{remark}

%══════════════════════════════════════════════════════════════════════════════
\section{The Generation-1 Electromagnetic Flip}
\label{P20:sec:gen1}
%══════════════════════════════════════════════════════════════════════════════

The static tower and the tangent-orientation scale give the up-type
quarks a heavier ordering at all generations; observation has $u<d$ at
generation~1 but $c>s$, $t>b$ at generations~2,3. This sign flip is not
a Coulomb self-energy: the self-energy $\sim+\alphaem Q^2$ makes the
higher-charge up heavier (wrong sign) and is $\sim0.5$~MeV (too small).

\begin{proposition}[Charge-linear generation-1 flip]
\label{P20:prop:gen1_flip}
A charge-linear coupling to a formation-scale field,
$\delta m=-Q\,\Phi$ with $\Phi\sim\alphaem\Lambda_{\mathrm{const}}
\approx1.6$~MeV, shifts up-type ($Q=+\tfrac23$) and down-type
($Q=-\tfrac13$) oppositely, with $\delta m_{\rm down}-\delta m_{\rm up}
=+\Phi>0$ (down heavier), of the observed magnitude $m_d-m_u=2.5$~MeV.
Because $\Phi$ is generation-independent, the shift is a fixed
$\sim$MeV offset: decisive at generation~1 ($\sim50\%$ of $m_u$) and
negligible at generations~2,3 (GeV masses). It therefore flips only
generation~1, reproducing the observed pattern.
\end{proposition}

\begin{remark}[Sign of the field]
\label{P20:rem:gen1_sign}
The sign of $\Phi$ (which member is raised) is fixed by observation, as
the down direction is fixed by the strange--muon anchor; this is the
external-field selection of Paper~XVIII's isospin conjecture
%\ref{P18:conj:isospin}
made quantitative. A first-principles derivation of $\Phi$ from the
density-feedback action is open.
\end{remark}

\begin{remark}[The flip is the neutron--proton QCD splitting]
\label{P20:rem:gen1_np}
The scale of $\Phi=m_d-m_u$ is not free: as the quark-mass isospin breaking it
drives the QCD contribution to the neutron--proton mass difference,
$(m_n-m_p)_{\rm QCD}\approx+2.5$~MeV (lattice), which makes the neutron the heavier
nucleon, while the electromagnetic self-energy supplies the opposite-sign
$(m_n-m_p)_{\rm QED}\approx-1$~MeV, the two summing to the observed $+1.29$~MeV. So
$\Phi$ is a precisely measured observable, not merely an ``order-MeV'' offset. Taking
the formation scale as the jam value $\Lambda\sim m_p^{\rm jam}/3\approx308$~MeV gives
$\Phi=\alphaem\Lambda\approx2.3$~MeV, $\sim90\%$ of the QCD splitting; the constituent
scale $237$~MeV gives $1.7$~MeV ($\sim70\%$). The magnitude is thus fixed by $\alphaem$
and the formation scale ($\sim250$--$310$~MeV) with no new coupling, the residual being
the sector's formation-scale ambiguity.
\end{remark}

%══════════════════════════════════════════════════════════════════════════════
\section{The Doublet Sum Rule (Empirical)}
\label{P20:sec:sumrule}
%══════════════════════════════════════════════════════════════════════════════

The down sector is anchored to the charged lepton; the up sector's
weak-isospin partner is the light neutrino, so the up sector has no
massive lepton anchor of its own (Paper~XVI~\cite{Paper16}). The
following relation supplies one indirectly, through the doublet.

\begin{conjecture}[Doublet mass-budget sum rule]
\label{P20:conj:sumrule}
For each confined generation the total quark mass of the weak doublet is
a fixed multiple of the charged-lepton mass,
\begin{equation}
  m_{\rm up}+m_{\rm down}=C\,m_\ell,\qquad C\approx13,
  \label{P20:eq:sumrule}
\end{equation}
holding across the two confined doublets:
$(m_u+m_d)/m_e=13.37$ and $(m_c+m_s)/m_\mu=12.90$ (a $3.5\%$ spread).
Using $C$ from one doublet to predict the other's up quark gives $m_c$
to $3.9\%$ and $m_u$ to $\sim11\%$ from the charged leptons and the
(lepton-anchored) down masses: $m_{\rm up}=C\,m_\ell-m_{\rm down}$. The
third doublet is excluded because the top is the bare electroweak-scale
transition (Paper~XVIII), not a confined member.
\end{conjecture}

\begin{remark}[What the sum rule resolves, and what it does not]
\label{P20:rem:sumrule_scope}
The relation ties the up sector to the \emph{same} charged lepton that
anchors the down sector, through the doublet budget rather than through
the up quark's own (light) neutrino partner: the charged-lepton mass is
the doublet budget, and the up quark takes the remainder after the
lepton-anchored down. This is the sense in which ``the charged-lepton
mass goes somewhere'' in the up sector. It is an empirical regularity
resting on two confined doublets; the magnitude of $C$ is not derived.
\end{remark}

\begin{remark}[Retracted numerical readings of $C$]
\label{P20:rem:C_retractions}
Several attractive readings of $C\approx13$ are recorded as \emph{not}
supported, to prevent their recurrence.
\emph{(i)~$C=4\pi+\tfrac13$.} The value $C\approx12.9$ is numerically
close to $4\pi+1/3=12.90$, with $1/3=h_{\SU(3)_1}$ the quark
$\SU(3)$ signature (equivalently the net doublet charge $+\tfrac13$).
But the $4\pi$ has no mechanism: in this framework the Hopf charge
enters mass \emph{exponentially} (the Chern--Simons spoke amplitude
$4\pi\,\Delta Q$ appears as $e^{4\pi}$ in the charged-lepton mass;
Paper~IV~\cite{Paper4}), so a mass ratio equal to the \emph{bare}
action $4\pi$ is a category error. The Chern--Simons $4\pi$ gives
$e^{4\pi}\approx3\times10^5$, the vacuum-polarisation $1/(4\pi)$ per
WZW level is an $\sim8\%$ correction, and the mass tower is $\vp^2$ per
level---none is a linear factor of $4\pi$.
\emph{(ii)~Solid-angle $4\pi$.} A direct computation of the director
$S^2$ solid angle (and of $J_2$, $J_4$) across the $\QH=2\to3$ step
(charge-verified hopfions) gives growth factors $1.2$--$1.8$, of order
the charge ratio $3/2$, not $4\pi$.
\emph{(iii)~$C=\Qgr^{\rm baryon}+\Qgr^{\rm lepton}=3+10$.} A clean
integer but with no mass-budget mechanism, and not distinguishable from
$4\pi+\tfrac13$ within the light-quark mass uncertainties.
The one principled fragment is the charge $1/3$; the magnitude of $C$
is an open, non-perturbative quantity.
\end{remark}

\begin{proposition}[The constituent scale from the $8\pi$ spoke, and the reduction of $C$]
\label{P20:prop:constituent}
The quark constituent scale follows from the same Chern--Simons and
condensate structure as the charged-lepton mass. With the baryon
condensate scale $\Lcond^{\rm baryon}=\TCMB(\pi^2/45)^{1/4}$ (the
$\Qgr=3$ member of $\rCMB=\Qgr\,\Lcond^4$; Paper~XII~\cite{Paper12}), the
tower factor $\vp^{2\Qgr}=\vp^6$, and the Chern--Simons spoke $e^{8\pi}$
(the $\Delta Q=3-1=2$ amplitude of the $\QH=3$ object above the $\QH=1$
base; Paper~V~\cite{Paper5}),
\begin{equation}
  m_{\rm const}=\Lcond^{\rm baryon}\,\vp^{6}\,e^{8\pi}\approx237~\mathrm{MeV},
  \label{P20:eq:constituent}
\end{equation}
matching the phenomenological constituent scale $215.5$~MeV to
$\sim10\%$---itself of the $\alphas$ residual size of
Section~\ref{P20:sec:residuals}. The ratio to the electron is
\begin{equation}
  \frac{m_{\rm const}}{m_e}
  =\Bigl(\tfrac{10}{3}\Bigr)^{1/4}\vp^{-14}\,e^{4\pi+3/400}\approx463,
  \label{P20:eq:const_over_me}
\end{equation}
the surviving factor being the \emph{extra} Chern--Simons spoke $e^{4\pi}$
of the quark ($\Delta Q=2$) over the charged lepton ($\Delta Q=1$).
Consequently the doublet sum-rule constant of
Conjecture~\ref{P20:conj:sumrule} is
\begin{equation}
  C=\frac{m_{\rm const}}{m_\ell}\,\vp^{-2\Delta n_q(g)},
  \label{P20:eq:C_reduction}
\end{equation}
with $\Delta n_q(g)$ the up-doublet tower level relative to the
constituent scale. Thus $C$ is fixed once $\Delta n_q(g)$---the
$E_6$-native tower level---is known: deriving $C$ and pinning the $E_6$
tower are the same problem.
\end{proposition}

\begin{remark}[$C$ is a residual, not a fundamental constant]
\label{P20:rem:C_residual}
Equation~\eqref{P20:eq:C_reduction} exhibits $C\approx13$ as a delicate
residual of the large factor $m_{\rm const}/m_\ell\approx463$ against the
tower suppression $\vp^{-2\Delta n_q}$, not a standalone number. This is
consistent with Remark~\ref{P20:rem:C_retractions}: the proximity to
$4\pi$ is coincidental, because $4\pi$ enters the underlying masses
\emph{exponentially} (as $e^{4\pi}$, $e^{8\pi}$) while $C$ is their
compensated ratio.
\end{remark}

%══════════════════════════════════════════════════════════════════════════════
\section{Residuals as the Dynamical Layer}
\label{P20:sec:residuals}
%══════════════════════════════════════════════════════════════════════════════

The results above are static/topological (group data, tangent geometry,
charge). They reproduce the observed masses to leading order and leave
residuals. We now show these residuals are not incidental but partition
by sector in a way that identifies them with a single dynamical
correction.

\begin{proposition}[The residual partition]
\label{P20:prop:residual_partition}
Classify the framework's quantitative residuals (Paper~VII~\cite{Paper7}
summary table, plus the quark-sector results above) by whether the
sector is a \emph{static saddle} (weakly-coupled, stable) or
\emph{dynamical} (strongly-coupled or loop-generated). Then:
\begin{itemize}[leftmargin=1.4em]
\item static-saddle sectors---the fine-structure constant
($5\times10^{-7}$), the electron mass ($1.3\times10^{-4}$), the Higgs
mass ($\sim10^{-3}$), $\sin^2\theta_W$ ($2.3\times10^{-3}$),
$m_\mu/m_\tau$ ($5\times10^{-3}$)---have residuals $\lesssim0.5\%$
(geometric mean $\approx0.011\%$);
\item dynamical sectors---the neutrino mass ($28\%$), the down-sector
running ($\sim20\%$), the isospin ratio ($21\%$), the up-type $u,c$
(factor)---have residuals $\gtrsim20\%$ (geometric mean $33\%$);
\item cosmological quantities, with mixed static and dynamical inputs,
lie between ($0.1$--$9\%$).
\end{itemize}
The two classes are separated by a factor $\sim40$ with no overlap.
\end{proposition}

\begin{remark}[Residual size tracks the sector coupling]
\label{P20:rem:coupling_scale}
The dynamical-sector residual scale, $\sim20$--$30\%$, is the
strong-coupling scale $\alphas\!\sim\!0.3$ at a few GeV, while the
static-sector residuals sit at $\sim\alphaem/\pi\sim2\times10^{-3}$.
The residual size therefore tracks the sector's dynamical coupling. The
interpretation is structural: the static/topological layer is the exact
leading order, and the residuals are one dynamical correction, largest
precisely where the object has no static saddle (the metastable quark)
or is loop-generated (the neutrino). This is not a collection of
unrelated misses.
\end{remark}

\begin{remark}[The perturbative strong layer is already built and does not close the residuals]
\label{P20:rem:perturbative}
Paper~V~\cite{Paper5} and Paper~XVI~\cite{Paper16} derive a
framework-internal $\alphas(\Lambda_{\rm UV})=C_F\sin(\pi/5)/(2\vp^{43/7})
\approx0.0204$ and run the quark masses by one-loop QCD with
flavour-threshold matching. This changes the predictions by at most a few
percent and leaves the pattern (one near-exact match at strange, five
undershoots, generation-growing) unchanged. The $\sim20$--$30\%$
residuals are therefore \emph{non-perturbative}: the metastable-quark
formation/confinement dynamics, consistent with $\alphas\sim0.3$ being
the regime where perturbation theory becomes order unity.
\end{remark}

\begin{remark}[Two distinct non-perturbative objects]
\label{P20:rem:two_objects}
The confinement string tension gives $\sigma R\approx278$~MeV and
$\sigma(3R_0+R)\approx925$~MeV$\,=98.6\%$ of $m_p$ (Paper~XV,
Paper~XVI): this is the constituent/\emph{hadron} scale, largely
accounted for, and is not a correction to \emph{current} quark masses.
The current-quark-mass residuals span six orders of magnitude in
absolute terms and are multiplicative, not a fixed $\sigma R$; they are
a distinct non-perturbative effect on the current-mass tower. The two
must not be conflated.
\end{remark}

\begin{table}[htbp]
\centering
\caption{Quark masses: framework prediction vs.\ measured (PDG,
$\overline{\mathrm{MS}}$; $b$ at $m_b$, light quarks at $2$~GeV, top pole).
Down-type use the $E_8$ bridge $m=m_\ell\,e^{n_q\pi/9}$ with the charged
leptons as input; their common $\sim\!1.2$ over-prediction is the
generation-universal RG running to the measurement scale. Up-type use the
$E_6$-native tower: the ratio $m_u/m_c=e^{-2\pi}$ is parameter-free
(Proposition~\ref{P20:prop:up_ratio}), and the absolute scale is set by the
completion (de-jamming) energy $(m_p-\sigma(3R_0+R))/3\approx4.5$~MeV
(Open Problem~O1), accurate to $\sim\!10\%$. The top sits at $v/\sqrt2$ by
incomplete formation---it decays before hadronising---while its full $T_{2,3}$
formation would be $m_\tau e^{15\pi/9}\approx334$~GeV. Every deviation is within
or comparable to the experimental current-mass uncertainty except at the
well-measured $c,b$, where it is the $\sim\!10$--$20\%$ shape/running residual.
This is agreement across five orders of magnitude, not a zero-parameter fit.}
\label{P20:tab:predvsmeas}
\begin{tabular}{clccc}
\hline
Quark & Route & Predicted & Measured (PDG) & Pred/Meas \\
\hline
$d$ & $E_8$ bridge, $m_e\,e^{7\pi/9}$   & $5.88$~MeV  & $4.67^{+0.48}_{-0.17}$~MeV & $1.26$ \\
$s$ & $E_8$ bridge, $m_\mu$             & $105.7$~MeV & $93.4^{+8.6}_{-3.4}$~MeV   & $1.13$ \\
$b$ & $E_8$ bridge, $m_\tau\,e^{\pi/3}$ & $5063$~MeV  & $4180\pm30$~MeV            & $1.21$ \\
$u$ & $E_6$-native, completion anchor   & $2.05$~MeV  & $2.16^{+0.49}_{-0.26}$~MeV & $0.95$ \\
$c$ & $E_6$-native, completion anchor   & $1098$~MeV  & $1270\pm20$~MeV            & $0.86$ \\
$t$ & incomplete formation, $v/\sqrt2$  & $174$~GeV   & $172.6\pm0.3$~GeV          & $1.01$ \\
\hline
\end{tabular}
\end{table}

%══════════════════════════════════════════════════════════════════════════════
\section{Open Problems}
\label{P20:sec:open}
%══════════════════════════════════════════════════════════════════════════════

\begin{enumerate}[leftmargin=2em,label=\textbf{O\arabic*.}]

\item \textbf{The absolute anchor of the $E_6$ tower.}
Proposition~\ref{P20:prop:constituent} derives the constituent scale
$\Lcond^{\rm baryon}\vp^6 e^{8\pi}$ and the ratio $m_{\rm const}/m_e$,
and reduces $C$ to the up-doublet tower level $\Delta n_q(g)$
(Eq.~\eqref{P20:eq:C_reduction}). The per-generation \emph{spacing} of that
tower is fixed with no free parameter by Proposition~\ref{P20:prop:up_ratio}
($m_u/m_c=e^{-2\pi}$ from the $\SU(2)_{k_g}$ T-matrix phases and the $E_6$
Coxeter exponents), so the $T_g$ half of the $n$--$T_g$ degeneracy (a single
mass fixes only $2n\ln\vp+2\pi\Qgr T_g$) is fixed by conformal data. What
remains is the single \emph{absolute} normalisation of the confined $u,c$
doublet---the $T_{2,2}\!\to\!T_{2,3}$ formation energy, \emph{not} the top:
the top decays before it can hadronise and sits at $v/\sqrt2$ by incomplete
formation, so it is not a rung of the confined tower
(Remark~\ref{P20:rem:up_ratio_data}). This is a formation-energy question in
the $(E_6)_1$ data.

The formation balance fixes this scale concretely. The confined tower anchors
not on the constituent (jammed hadron) scale $\Lcond^{\rm baryon}\vp^6 e^{8\pi}
\approx237$~MeV but on the \emph{completion} energy---the de-jamming/localisation
piece $m_p-\sigma(3R_0+R)\approx13.5$~MeV released at hadronisation, i.e.\
$\approx4.5$~MeV per quark, matching the fitted $u,c$-tower anchor ($\approx5$~MeV)
to $\sim10\%$.

A residual of this order is \emph{expected}, and distinguishes the sector. The
charged lepton ($\QH=2$) is a free particle whose isolated configuration exists
and fixes $m_e$ to $0.013\%$ (Paper~IV~\cite{Paper4}). The quark ($\QH=3$) has
no stable isolated
configuration: it holds its shape only in the jam, through the shared lobes and
crossings of the confined network, so its mass is intrinsically the \emph{shared}
value and never the isolated topological ideal from which the WZW tower is
computed. The per-generation shape residual ($\sim5$--$10\%$: the $u/c$ ratio and
the $d$--$s$ split at a common scale, distinct from the generation-universal RG
running of the down sector) is therefore the irreducible shared-jam deviation of a
confined quark from its free-particle ideal---a deviation the free leptons,
realising their ideal, do not carry. It currently lies within the experimental
current-mass uncertainties ($\sim{+23}/{-12}\%$ for $m_u$, $\sim10\%$ for $m_d$),
so it stands as a \emph{prediction}---confined-quark towers should carry a
shared-jam deviation absent from the leptons---rather than a resolved discrepancy.
Computing the completion energy and this deviation from first principles requires
the jammed (topology-protected) $\QH=3$ configuration: the isolated $\QH=3$ trefoil
has no stable saddle, so a topology-protected treatment (not unconstrained gradient
flow) is a prerequisite.

\item \textbf{Derivation of the generation-1 field $\Phi$.}
Proposition~\ref{P20:prop:gen1_flip} fixes the charge-linear flip
$\Phi=\alphaem\Lambda$ to the formation scale, and
Remark~\ref{P20:rem:gen1_np} identifies $\Phi$ with the QCD part of the
neutron--proton mass difference ($\approx2.5$~MeV), reached to
$\sim70$--$90\%$ across the constituent-to-jam scale range. What remains
is a first-principles $\Phi$ from the density-feedback action---both its
magnitude to better than the formation-scale ambiguity, and its sign
(currently set by the strange--muon anchor)---and the connection to the
Paper~XVIII isospin field.

\item \textbf{First-principles $\Mdn/\Mup$ residual.}
Proposition~\ref{P20:prop:isospin_scale} closes to $\approx1.2$. This is
\emph{not} the down-sector running---which cancels in a same-scale
up/down ratio---so the two are distinct corrections of coincidentally
similar size. What remains open is whether the tangent-fraction
functional is derivable from the density-feedback action rather than
motivated.

\end{enumerate}

%══════════════════════════════════════════════════════════════════════════════
\section{Summary}
\label{P20:sec:summary}
%══════════════════════════════════════════════════════════════════════════════

The $\QH=3$ quark sector separates into a static/topological leading
order and a dynamical correction. The static layer comprises: the
$E_6$-native generation structure with the $E_8$ charged-lepton bridge
(Section~\ref{P20:sec:groups}), giving the up/down steepness as the WZW
phase-unit ratio $9/4$; the isospin mass-scale from the trefoil tangent
orientation (Section~\ref{P20:sec:isospin}); the fractional charge from
the colour $\Zthree$ center (Section~\ref{P20:sec:charge}), derived
rather than assigned; the charge-linear generation-1 flip
(Section~\ref{P20:sec:gen1}); the derived constituent scale
$m_{\rm const}=\Lcond^{\rm baryon}\vp^6 e^{8\pi}$
(Proposition~\ref{P20:prop:constituent}), which reduces the sum-rule
constant to the $E_6$ tower level; and the empirical
doublet sum rule (Section~\ref{P20:sec:sumrule}), whose one derived
fragment is the $\SU(3)$ charge and whose magnitude is open.

The organising result (Section~\ref{P20:sec:residuals}) is that the
framework's residuals partition by sector: static-saddle sectors are
predicted to $\lesssim0.5\%$, dynamical sectors carry
$\sim\!\alphas\sim20$--$30\%$ residuals, and the perturbative strong
layer (a framework-internal $\alphas$, already built) does not close
them. The residuals are the non-perturbative dynamical layer, forced
into view in the quark sector because the quark alone among the massive
fermions has no static saddle. The two un-developed handles---the
Chern--Simons edge/Hall dynamics of fractional charge and the
metastable-quark formation energy---are the natural next targets, and
are the content this framework must supply beyond its static skeleton.

%══════════════════════════════════════════════════════════════════════════════
\begin{thebibliography}{99}

\bibitem{Paper1}
F.~Manfredi,
\textit{The Density-Feedback Faddeev--Niemi Hopfion:
  Exact Algebraic Predictions for Standard-Model Parameters},
Zenodo (2026).
\href{https://doi.org/10.5281/zenodo.19342027}{doi:10.5281/zenodo.19342027}

\bibitem{Paper4}
F.~Manfredi,
\textit{The Hopf Spoke, WZW Fermion Mass Renormalization,
  and Three- and Four-Term Formulas for the Fine Structure Constant},
Zenodo (2026).
\href{https://doi.org/10.5281/zenodo.19504729}{doi:10.5281/zenodo.19504729}

\bibitem{Paper5}
F.~Manfredi,
\textit{Colour Confinement, the QCD Scale, and Hadronic Structure
  from the Density-Feedback Hopfion},
Zenodo (2026).
\href{https://doi.org/10.5281/zenodo.19638857}{doi:10.5281/zenodo.19638857}

\bibitem{Paper6}
F.~Manfredi,
\textit{The Golden-Ratio Tower: Fermion Scale Hierarchy from the
  Hopfion RG Fixed Point},
Zenodo (2026).
\href{https://doi.org/10.5281/zenodo.19646945}{doi:10.5281/zenodo.19646945}

\bibitem{Paper7}
F.~Manfredi,
\textit{Emergent Gravity, Inflation, Dark Energy, and the Weak
  Equivalence Principle from the Density-Feedback Hopfion Condensate},
Zenodo (2026).
\href{https://doi.org/10.5281/zenodo.19646953}{doi:10.5281/zenodo.19646953}

\bibitem{Paper12}
F.~Manfredi,
\textit{The Profile Normalisation Conjecture ---
  Proof via CMB Energy Identity and Two-Route Equivalence},
Zenodo (2026).
\href{https://doi.org/10.5281/zenodo.20210460}{doi:10.5281/zenodo.20210460}

\bibitem{Paper15}
F.~Manfredi,
\textit{The $Q_H=3$ Sector of the Density-Feedback Hopfion},
Zenodo (2026).
\href{https://doi.org/10.5281/zenodo.20691001}{doi:10.5281/zenodo.20691001}

\bibitem{Paper16}
F.~Manfredi,
\textit{Quark Generation Masses in the Density-Feedback Hopfion:
A Survey of Ruled-Out Mechanisms and the $Q_H=3$ Gradient-Flow Landscape},
Preprint (2026).
\href{https://doi.org/10.5281/zenodo.21012650}{doi:10.5281/zenodo.21012650}

\bibitem{Paper17} F.~Manfredi,
\textit{The $Q_H=1$ Sector of the
Density-Feedback Hopfion: Topology, Group Structure, Colour Exclusion},
Zenodo (2026).
\href{https://doi.org/10.5281/zenodo.21013412}{doi:10.5281/zenodo.21013412}

\bibitem{Paper18}
F.~Manfredi,
\textit{Preimage Topology and Confinement: The $Q_H=3$ Sector of the Density-Feedback Hopfion},
Zenodo (2026) - DRAFT.
\href{https://doi.org/10.5281/zenodo.21047750}{doi:10.5281/zenodo.21047750}

\bibitem{Paper19}
F.~Manfredi,
\textit{Dynamical Origin of Quark Masses in the Density-Feedback Hopfion},
Zenodo (2026).
\href{https://doi.org/10.5281/zenodo.21225452}{doi:10.5281/zenodo.21225452}

\bibitem{Nayak2008}
C.~Nayak, S.~H.~Simon, A.~Stern, M.~Freedman, S.~Das~Sarma,
\textit{Non-Abelian anyons and topological quantum computation},
Rev.\ Mod.\ Phys.\ \textbf{80}, 1083--1159 (2008).
\href{https://arxiv.org/abs/0707.1889}{arXiv:0707.1889}

\bibitem{Everett2009}
L.~L.~Everett, A.~J.~Stuart,
\textit{Icosahedral ($A_5$) family symmetry and the golden ratio
prediction for solar neutrino mixing},
Phys.\ Rev.\ D \textbf{79}, 085005 (2009).
\href{https://arxiv.org/abs/0812.1057}{arXiv:0812.1057}

\bibitem{Feruglio2017}
F.~Feruglio,
\textit{Are neutrino masses modular forms?},
in \textit{From My Vast Repertoire}, World Scientific (2019), pp.\ 227--266.
\href{https://arxiv.org/abs/1706.08749}{arXiv:1706.08749}

\bibitem{ModFlavHierarchy2025}
M.-C.~Chen et al.,
\textit{Modular flavor symmetries and fermion mass hierarchies},
(2025).
\href{https://arxiv.org/abs/2506.23343}{arXiv:2506.23343}

\bibitem{FracCharge2024}
R.~Alonso et al.,
\textit{Fractional-charge hadrons and leptons to tell the Standard Model
group apart},
Phys.\ Lett.\ B \textbf{863}, 139354 (2025).
\href{https://arxiv.org/abs/2404.03438}{arXiv:2404.03438}

\end{thebibliography}

\end{document}

